\section*{第 \ref{ch:g12:condprob} 章——条件概率与独立性}

\begin{solution}{exo:g12:condprob:1}
共有 $12$ 张人头牌，其中 $4$ 张是国王：
$\pcond{\text{人头牌}}{\text{国王}} = \frac{4}{12} = \frac13$。

\omterm{def:g12:condprob:indep}{独立}性：$\P(\text{国王} \cap \text{红心}) = \frac{1}{52}$（红心
国王），且
$\P(\text{国王})\,\P(\text{红心}) = \frac{4}{52}\times\frac{13}{52}
= \frac{1}{52}$。两者相等，因此两\omterm{def:g11:prob:model}{事件}\omterm{def:g12:condprob:indep}{独立}。
\end{solution}

\begin{solution}{exo:g12:condprob:2}
\emph{1.} 第一步分支：红球 $\frac58$，蓝球 $\frac38$；第二步分支
（不放回）：取红后，红 $\frac47$ / 蓝 $\frac37$；
取蓝后，红 $\frac57$ / 蓝 $\frac27$。

\emph{2.} $\P(RR) = \frac58 \times \frac47 = \frac{5}{14}$。由全
\omterm{def:g10:proba:distribution}{概率}公式，
\[
\P(\text{第二次为红}) = \frac58\cdot\frac47 + \frac38\cdot\frac57
= \frac{20 + 15}{56} = \frac58 .
\]
（与第一次抽取的结果相同——由对称性，第二个球就是瓮中
一个均匀随机抽取的球。）
\end{solution}

\begin{solution}{exo:g12:condprob:3}
$\P(A) = \frac{6}{36} = \frac16$（点数和为 $7$ 的数对共六组），
$\P(B) = \frac16$，且 $A \cap B = \{(3,4)\}$ 的\omterm{def:g10:proba:distribution}{概率}为
$\frac{1}{36} = \P(A)\P(B)$，故 $A$ 与 $B$ \omterm{def:g12:condprob:indep}{独立}。（和为 $7$ 的
情形很特殊：无论第一颗骰子掷出什么，第二颗骰子恰有一个点数使两数之和为该值。）

$\P(C) = \frac{5}{36}$，且 $B \cap C = \{(3,3)\}$：
$\P(B \cap C) = \frac{1}{36} \neq \frac16 \times \frac{5}{36}
= \frac{5}{216}$。两\omterm{def:g11:prob:model}{事件}不\omterm{def:g12:condprob:indep}{独立}。
\end{solution}

\begin{solution}{exo:g12:condprob:4}
\emph{1.} 按机器划分，用全\omterm{def:g10:proba:distribution}{概率}公式：
\[
\P(D) = 0.5\times0.01 + 0.3\times0.02 + 0.2\times0.04
= 0.005 + 0.006 + 0.008 = 0.019 = 1.9\% .
\]

\emph{2.} 由贝叶斯公式：
$\pcond{D}{M_3} = \dfrac{0.008}{0.019} = \dfrac{8}{19} \approx 0.42$。这说明
仅生产五分之一产量的机器，
却贡献了超过 $40\%$ 的次品。
\end{solution}

\begin{solution}{exo:g12:condprob:5}
条件是
\[
\frac{0.99\,p}{0.99\,p + 0.05\,(1-p)} \geq 0.9 .
\]
分母为正，故上式可化为
$0.99p \geq 0.891p + 0.045 - 0.045p$，\emph{即} $0.144\,p \geq 0.045$，
\emph{即}
\[
p \geq \frac{0.045}{0.144} = 0.3125 .
\]
仅靠这一检验要达到 $90\%$ 的可信度，只有当该病已
使超过 $31\%$ 的受检人群患病——这正是罕见病大规模筛查
需要进行确认检验的原因。
\end{solution}

\begin{solution}{exo:g12:condprob:6}
\emph{1.} 满足 $E$ 的序列有三个（HHT、HTH、THH），每个的\omterm{def:g10:proba:distribution}{概率}
均为 $p^2(1-p)$，由\omterm{def:g12:condprob:indep}{独立}性得 $\P(E) = 3p^2(1 - p)$。

\emph{2.} 令 $f(p) = 3p^2 - 3p^3$，则 $f'(p) = 6p - 9p^2 = 3p(2 - 3p)$，
在 $\intoo{0}{1}$ 上先正后负，故\omterm{def:g10:functions:extrema}{最大值}在 $p = \frac23$ 处取得，此时
$\P(E) = 3\cdot\frac49\cdot\frac13 = \frac49$。
\end{solution}

\begin{solution}{exo:g12:condprob:7}
设你选择了 1 号门，主持人打开了 3 号门（\omterm{def:g11:prob:model}{事件} $H_3$）。记
$B_i$ 为“奖品在 $i$ 号门后”，$\P(B_i) = \frac13$，且
\[
\pcond{B_1}{H_3} = \frac12, \qquad
\pcond{B_2}{H_3} = 1, \qquad
\pcond{B_3}{H_3} = 0
\]
（若奖品在你所选的门后，主持人便在 2 号门与 3 号门之间
随机选择；若奖品在 2 号门后，他则被迫打开 3 号门）。由贝叶斯公式：
\[
\pcond{H_3}{B_1}
= \frac{\frac13\cdot\frac12}{\frac13\cdot\frac12 + \frac13\cdot1 + 0}
= \frac{1/6}{1/2} = \frac13,
\qquad
\pcond{H_3}{B_2} = \frac{\frac13}{\frac12} = \frac23 .
\]
不换门获胜的\omterm{def:g10:proba:distribution}{概率}为 $\frac13$，换门获胜的\omterm{def:g10:proba:distribution}{概率}为
$\frac23$。
\end{solution}

\begin{solution}{exo:g12:condprob:8}
\emph{1.} 多数表决出错，当且仅当三个副本中至少有两个
被翻转，\omterm{def:g10:proba:distribution}{概率}为：
\[
3\varepsilon^2(1-\varepsilon) + \varepsilon^3
= 3(0.01)(0.9) + 0.001 = 0.028,
\]
远小于 $\varepsilon = 0.1$：重复编码有效。

\emph{2.} 若发送 $1$（即 $111$），收到 $101$ 恰需一次翻转，
其\omterm{def:g10:proba:distribution}{概率}为 $\varepsilon(1-\varepsilon)^2 = 0.081$。若发送 $0$
（即 $000$），则需两次翻转：$\varepsilon^2(1-\varepsilon) = 0.009$。
设先验\omterm{def:g10:proba:distribution}{概率}相等，由贝叶斯公式：
\[
\pcond{101}{\,\text{发送 }1} = \frac{0.081}{0.081 + 0.009} = 0.9 .
\]
多数译码（$101 \mapsto 1$）确实是更可能的猜测。
\end{solution}

\begin{solution}{pb:g12:condprob:1}
\textbf{1.} $\P(\text{红}) = \frac12 \cdot \frac23 + \frac12
\cdot \frac14 = \frac13 + \frac18 = \frac{11}{24}$。

\textbf{2.} $\P(\text{I} \mid \text{红球}) =
\dfrac{\frac12 \cdot \frac23}{\frac{11}{24}} =
\dfrac{1/3}{11/24} = \dfrac{8}{11}$。

\textbf{3.} $\P(\text{偶数} \cap \leq 4) = \P(\{2, 4\}) =
\frac13$，而 $\P(\text{偶数})\P(\leq 4) = \frac12 \cdot
\frac23 = \frac13$，\omterm{def:g12:condprob:indep}{独立}（出人意料！）。对于
$\leq 3$：$\P(\{2\}) = \frac16$，
对比 $\frac12 \cdot \frac12 = \frac14$：不\omterm{def:g12:condprob:indep}{独立}。

\textbf{4.} \omterm{def:g12:condprob:indep}{独立}\omterm{def:g11:prob:model}{事件}的
\omterm{def:g10:proba:operations}{对立事件}
仍然\omterm{def:g12:condprob:indep}{独立}，因此
有
\[
\P(\text{一个都不发生}) = 0.1 \times 0.1 = 0.01,
\qquad
\P(\text{至少一个发生}) = 0.99 .
\]

\textbf{5.} $\frac{4}{52} \times \frac{3}{51} =
\frac{1}{221}$。

\textbf{6.} 患病者：$100$；健康者：$99\,900$。阳性结果：
$99$ 个真阳性（$100$ 的 99\,\%），$4\,995$ 个假阳性（5\,\% 的
$99\,900$），共 $5\,094$ 个阳性结果。

\textbf{7.} $\P(\text{患病} \mid +) = \frac{99}{5\,094}
\approx 1.9\,\%$：在检测判为阳性的人中，
真正患病者不足五十分之一。

\textbf{8.} $\dfrac{0.001 \times 0.99}{0.001 \times 0.99 +
0.999 \times 0.05} = \dfrac{0.00099}{0.05094} \approx 0.019$：
计数与公式结果一致。

\textbf{9.} 混淆了 $\P(+ \mid \text{患病}) = 99\,\%$
（检测宣称的准确率）与
$\P(\text{患病} \mid +)$（患者真正关心的问题）。
在基础\omterm{def:g10:proba:distribution}{概率}为 $\frac{1}{1000}$ 的情形下，患病者极为罕见，
即便庞大的健康人群中仅 $5\,\%$ 的误报，
也会以五十比一的倍数淹没真阳性。

\textbf{10.} $\dfrac{0.019 \times 0.99}{0.019 \times 0.99 +
0.981 \times 0.05} \approx 28\,\%$。每次\omterm{def:g12:condprob:indep}{独立}的阳性结果
都会使几率乘以同一因子：证据不断累积，
第三次阳性将使\omterm{def:g10:proba:distribution}{概率}超过 $90\,\%$——这正是
需要复核性检测的原因。

\textbf{11.} $\dfrac{0.05 \times 0.99}{0.05 \times 0.99 +
0.95 \times 0.05} \approx 51\,\%$：在高风险人群中，
同一检测便具有参考价值。教训：应在基础\omterm{def:g10:proba:distribution}{概率}
可观的人群中筛查；对罕见病大规模筛查，
会制造出成千上万的误报。

\textbf{12.} 无辜者的预期吻合人数：
$10^7 \times 10^{-6} = 10$。在约 $11$ 名吻合者中
（$10$ 名无辜、$1$ 名有罪——若真凶在该城），
而由\emph{仅凭吻合}找到的嫌疑人，
是无辜者的\omterm{def:g10:proba:distribution}{概率}约为 $\frac{10}{11} \approx
91\,\%$。

\textbf{13.} 检察官引用的是
$\P(\text{吻合} \mid \text{无辜}) = 10^{-6}$；法庭
需要的是 $\P(\text{无辜} \mid \text{吻合})$——这正是
第 9 题中的混淆在宣誓之下重演。（在使这些谬误广为人知的真实案件中，
另一个失误使其更为严重：
见第 14 题。）

\textbf{14.} 直接平方：$\left(\frac{1}{8500}\right)^2
\approx \frac{1}{72 \text{ 百万}}$。计入第二次惨剧
在已知第一次发生的条件下可能性高出五倍：
$\frac{1}{8500} \times \frac{5}{8500} \approx
\frac{1}{14.5 \text{ 百万}}$——增大了五倍，却仍然忽略了
\emph{备择假设}同样须由贝叶斯公式加以权衡。
教训：\omterm{def:g10:proba:distribution}{概率}相乘
需要\omterm{def:g12:condprob:indep}{独立}性，而共同成因（如遗传、
环境）会将其破坏。

\textbf{15.} 贝叶斯公式将
\emph{全部}证据的似然比相乘，合成为一个后验\omterm{def:g10:proba:distribution}{概率}：吻合、
不在场证明、纤维证据，每一项都在改变几率。陪审员的
准则是：\emph{问一问证据在两种假设下各有多大可能，
切勿把某一个\omterm{def:g12:condprob:cond}{条件概率}误当作
其逆}。

\textbf{16.} $\dfrac{0.4 \times 0.15}{0.4 \times 0.15 +
0.6 \times 0.005} = \dfrac{0.060}{0.063} \approx 95\,\%$
垃圾邮件：一个词几乎就能定论。

\textbf{17.} 五份副本在 $3$、$4$ 或 $5$ 位翻转时出错：
$\binom53 \varepsilon^3(1-\varepsilon)^2 + \binom54
\varepsilon^4 (1 - \varepsilon) + \varepsilon^5 \approx
0.0086$——三份副本则为 $0.028$：
重复编码以带宽为代价换取可靠性。

\textbf{18.} 偶数次翻转：$\binom40 0.9^4 + \binom42 0.9^2
0.1^2 + \binom44 0.1^4 = 0.6561 + 0.0486 + 0.0001 = 0.7048$；
而 $\frac{1 + 0.8^4}{2} = \frac{1 + 0.4096}{2} = 0.7048$：
这一技巧公式追踪偏差 $(1 - 2\varepsilon)$
在各级中继中的传递情况。

\textbf{19.} 垃圾邮件用词常成群结伙出现——含有“彩票”的邮件
很可能也含有“中奖”，因此两者的联合证据
比乘积假设的情形弱。过滤器
仍然有效，因为这种夸大通常推向
\emph{正确的方向}：分类只需判断
结论位于 $50\,\%$ 的正确一侧，而非后验\omterm{def:g10:proba:distribution}{概率}本身正确。

\textbf{20.} 诊所：罕见病把强检测
变成弱结论——先由基础\omterm{def:g10:proba:distribution}{概率}更新，
再逐步累积。法庭：大城市中罕见的吻合
指控的是这座城市，而不是那个人——似然不等于定罪。
收件箱：许多弱线索相乘，便能分拣邮件——即便
只是朴素地运用，也能奏效。三个陷阱：忘记先验\omterm{def:g10:proba:distribution}{概率}，便会受骗；
把相依误当作\omterm{def:g12:condprob:indep}{独立}，便会大受其骗；零散地读取证据，
便永远不清楚自己究竟知道什么。
\end{solution}
